\documentclass[10pt,a4paper]{article}
\usepackage{amsmath,amssymb,amsthm,graphicx}
\title{Small circuits for truncations of low-degree rational functions}
\author{Pranjal Dutta, Gorav Jindal, Anurag Pandey, Amit Sinhababu}
\date{doi:10.4230/LIPIcs.CCC.2021.25; CC BY4.0}
\begin{document}
\maketitle
\noindent\textit{Editorial note.} Original authors: Pranjal Dutta, Gorav Jindal, Anurag Pandey, Amit Sinhababu. Official CCC2021 article25, DOI10.4230/LIPIcs.CCC.2021.25.\par\medskip
\noindent\textit{Editorial note.} License: https://creativecommons.org/licenses/by/4.0/\par\medskip
\noindent\textit{Editorial note.} Faithful editable TeX transcription of formal Theorem26 and complete author Section5.1 proof, full derivative Lemma63 and valuation Lemma65 proofs, exact partial-fraction/Leibniz facts and preceding complete division-elimination proofs as necessary context. Original field/degree/formal-series conditions, expressions and complete published bibliography retained without mathematical correction or new proofs. Only official CC BY4.0 version used.\par\medskip
\noindent\textit{Allowed standard prerequisites (editorial).}\begin{itemize}
\item Coefficient extraction and homogenization. Exact author Lemmas61/62 credited to [43] and [45] retained.
\item Polynomial Identity Lemma. Exact author Lemma64[38,15,51,41] retained.
\item Partial fraction decomposition and Leibniz rule. Exact author folklore Lemma24 and standard Lemma66 retained.
\end{itemize}\medskip
\section*{Author preliminaries}
We denote $\mathbf x=(x_1,\ldots,x_n)$. $[n]$ denotes $\{1,\ldots,n\}$. For a polynomial $f\in\mathbb F[\mathbf x]$, we denote up to degree-$d$ part as $\mathrm{Hom}_{\le d}f$ and $|f|_0$ as the sparsity or the number of monomials in $f$. For a differentiable function $f(x)$, we denote $f^{(k)}(x):=d^kf/dx^k$, as the $k$-th derivative of $f$. We also recall the definition of gcd of two polynomials $f,g$ in the ring $\mathbb F[\mathbf x]$: $\gcd(f,g)=:h\iff h\mid f,h\mid g,\text{ and }h'\mid f,g\implies h'\mid h$. It is unique up to constant multiples.
\paragraph{Field.} We denote the underlying field as $\mathbb F$ and assume that it is algebraically closed. All our results hold when the characteristic is large or not algebraically closed, as we can go to polynomial extensions and work with it.
\paragraph{div and mod operations.} For polynomials $f$ and $g\in\mathbb F[\mathbf x]$, if $f=g\cdot h+r$, where $h,r\in\mathbb F[\mathbf x]$ such that $\deg(r)<\deg(g)$, then $h$ is called the quotient, denoted $f\mathbin{\mathrm{div}}g$, and $r$ is called the remainder, denoted $f\mathbin{\mathrm{mod}}g$. Operation $\mathrm{mod}$ may not be well-defined in the multivariate settings, however, if one assumes $g$ to be monic in a variable say $x_n$, it is always well-defined (by thinking $g$ to be a univariate in $x_n$). A polynomial $g$ is monic in $x_n$ if the leading coefficient (the nonzero coefficient of highest degree) of $x_n$ is a non-zero constant in $\mathbb F$. Of course, if $g\mid f$, then $f\mathbin{\mathrm{div}}g=h$ and $f\mathbin{\mathrm{mod}}g=0$, irrespective of monic-ness.
\paragraph{Circuit complexity.} For a polynomial $f\in\mathbb F[\mathbf x]$, we define the complexity $L_{\mathbb F}(f)$ of $f$ to be the length of the smallest division-free arithmetic circuit (with only $\{+,-,\times\}$ gates) computing $f$. We also define the complexity $\tau_{\mathbb F}(f)$ of $f$ to be the length of the smallest division and constant-free arithmetic circuit computing $f$ (all the constants are made from 1), for formal definition see Section 5.2 of the original paper. We will remove subscript $\mathbb F$ when the underlying field is clear from the context.
\paragraph{Border complexity context (author Section 3.3).} The notion of border (equivalently, approximative) complexity is important in computer science. This concept popped up from early works on matrix multiplication and border rank of tensors, see [10]). Whether approximation of polynomials provides any additional computational power is a natural question which fundamentally motivated the foundation of Geometric Complexity theory (GCT). The notion of border complexity can be motivated through two ways: topological and algebraic, and both the perspectives are known to be equivalent [1]. For further details, we refer to [20,36]. In this paper, we only work with algebraic approximation upper bounds. In the algebraic definition, one can talk about the convergence $\epsilon\to0$. Here, one can see $\epsilon$ as a formal variable and $\mathbb F(\epsilon)$ as the function field. For an algebraic complexity class $\mathcal C$, the approximation is defined as follows [7,Definition2.1].
\paragraph{Definition 20 (Approximative closure of a class [7]).} Let $\mathcal C$ be an algebraic complexity class over field $\mathbb F$. A family $(f_n)$ of polynomials from $\mathbb F[\mathbf x]$ is in the class $\overline{\mathcal C}(\mathbb F)$ if there are polynomials $f_{n;i}$ and a function $t:\mathbb N\mapsto\mathbb N$ such that $g_n$ is in the class $\mathcal C$ over the field $\mathbb F(\epsilon)$ with $g_n(\mathbf x)=f_n(\mathbf x)+\epsilon f_{n,1}(\mathbf x)+\epsilon^2f_{n,2}(\mathbf x)+\cdots+\epsilon^{t(n)}f_{n,t(n)}(\mathbf x)$.
\paragraph{Definition 21 ([8,Defn.3.1]).} Let $f\in\mathbb F[\mathbf x]$. The border complexity $\underline L(f)$ is the smallest number $r$, such that there exists $F$ in $\mathbb F(\epsilon)[\mathbf x]$ satisfying $F|_{\epsilon=0}=f$ and $L_{\mathbb F(\epsilon)}(F)\le r$.
Note that the circuit of $F$ may be using $1/\epsilon$ in an intermediate step. So, we cannot merely assign $\epsilon=0$ and get a $\epsilon$-free circuit. Also, the $\epsilon$-degree can be exponential in its size (and thus cannot be interpolated), see [8,Theorem5.7]). Thus, potentially $\underline L(f)$ can be significantly smaller than $L(f)$. The above definition can be used to define closures of complexity class, e.g., $\overline{\mathrm{VP}}$. In this case, one can assume wlog that the degrees of $g_n$ and $f_{n,i}$ are $\mathrm{poly}(n)$. It is known to be closed under factoring [8,Theorem4.1]. However, the usual method of Hensel-lifting does not work when the given circuit class computes polynomials of super-polynomial degree. Also, Strassen's method would have a dependency on the degree of the final polynomial. However, we can prove Theorem 18 analogously, in the border sense.
\paragraph{Lemma 61 (Coefficient-Extraction).} Coefficient-extraction in arithmetic circuits is easy using interpolation, see the folklore lemma below, for a proof see [43]. Let $L(f)=s$ with $f\in\mathbb F[\mathbf x]$ and $f=\sum_{0\le i\le d}f_ix_n^i$ with $f_i\in\mathbb F[x_1,\ldots,x_{n-1}]$. Then there is a circuit $C$ of size $O(sd^2)$ computing $f_0,f_1,\ldots,f_d$.
\paragraph{Lemma 62 (Homogenization).} The next lemma is a homogenization trick, used in [45]. For a proof, see [43,Theorem2.2]. If $f$ has an arithmetic circuit of size $s$, then for any $d$, there is a circuit of size $O(sd^2)$ computing $\mathrm{Hom}_{\le d}f$.
\paragraph{Lemma 64 (Polynomial Identity Lemma).} Let $p(\mathbf x)$ be an $n$-variate nonzero polynomial of degree $d$. Let $S\subseteq\mathbb F$ be a finite set. Then $\Pr_{\alpha\sim S^n}[p(\alpha)=0]\le d/|S|$. Here, $\alpha\in S^n$ is picked independently and uniformly at random. This is the standard Polynomial Identity Lemma [38,15,51,41].

\section*{Formal author division-elimination main results}
\paragraph{Theorem 18 (Division elimination for low-degree divisor).} Let the polynomials $f,g,h\in\mathbb F[\mathbf x]$ such that $f=g/h$, with $L(g)=s_1$, $L(h)=s_2$, and $\deg(h)=d$. Then, $L(f)\le O((s_1+s_2)d^2)$.
\paragraph{Theorem 22 (Division elimination in border complexity).} Let $f,g,h\in\mathbb F[\mathbf x]$, such that $f=g/h$, with $\underline L(g)=s_1$, $\underline L(h)=s_2$, and $\deg(h)=d$. Then, $\underline L(f)\le O(s_1d^2+s_2d^4)$.

\section*{Complete author low-degree division tools}
\paragraph{Theorem 13.} Let $g,h$ be polynomials in $\mathbb F[x]$. If $L(g)=s$ and $\deg(h)=d$, then both $L(g\mathbin{\mathrm{div}}h)$ and $L(g\mathbin{\mathrm{mod}}h)$ have complexity $O(sd)$.
\begin{proof}Suppose $C$ is a circuit of size $s$ which computes $g$. We split every gate $\Phi$ in $C$ into two gates $\Phi_1$ and $\Phi_2$, to make a new circuit $C'$, which computes both $g\mathbin{\mathrm{div}}h$ and $g\mathbin{\mathrm{mod}}h$. If $\Phi$ is computing some polynomial $\phi$ in $C$, then $\Phi_1$ computes the polynomial $\phi\mathbin{\mathrm{mod}}h$ and $\Phi_2$ computes the polynomial $\phi\mathbin{\mathrm{div}}h$.
The proof is inductive and traverses from bottom to the top. The base case is trivial. At some step, say that we are at a gate $\Phi$. The children gate of $\phi$ are computing polynomials $\alpha$ and $\beta$. Let $\alpha=q_1h+r_1$, $\beta=q_2h+r_2$ and $\phi=qh+r$, where the degrees of $r,r_1,r_2$ are smaller than $d$. So in the new circuit $C'$, we have already computed $r_1,q_1,r_2,q_2$. If $\Phi$ is a $\pm$ gate then it is clear that $r=r_1\pm r_2$ and $q=q_1\pm q_2$. If $\Phi$ is a $\times$ gate then we have:
\begin{align*}r&=(r_1r_2)\mathbin{\mathrm{mod}}h,\\q&=q_1q_2h+r_1q_2+r_2q_1+(r_1r_2)\mathbin{\mathrm{div}}h.\end{align*}
We know that $r$ is a polynomial of degree at most $d-1$. Since $\deg(r_1r_2)\le2d-2$, we get that $\deg((r_1r_2)\mathbin{\mathrm{div}}h)\le d-2$. Therefore, we trivially have that: $L(r)=O(d)$ and $L((r_1r_2)\mathbin{\mathrm{div}}h))=O(d)$. Hence we can compute $r,q$ using additional $O(d)$ many gates. Thus, $C'$ has at most $O(sd)$ many gates. Hence $L(g\mathbin{\mathrm{div}}h)=O(sd)$ (same for $g\mathbin{\mathrm{mod}}h$).\end{proof}
\paragraph{Corollary 14.} For $f,g,h\in\mathbb F[x]$, if $f=g/h$ with $L(g)=s$ and $\deg(h)=d$ then $L(f)=O(sd)$.
\paragraph{Remark 15.} The polynomial $f_d:=1+\cdots+x^d=(x^{d+1}-1)/(x-1)$ has $O(\log d)$ size circuit. This can also be shown via a recursive computation argument. Can we expect both $\mathrm{div}$ and $\mathrm{mod}$ to have $\mathrm{poly}(s,\log d)$-size circuits? We show that it is highly unlikely unless factoring is easy, see Theorem 54 in the original paper for details.
\paragraph{Lemma 16 (same as Appendix E Lemma 70).} Suppose $g=\sum_{i\le d_1}g_ix^i$ and $h=x^{d_2}+\sum_{i<d_2}h_ix^i$, in $\mathbb F[x]$. Suppose $g=hq+r$, with $r=\sum_{i<d_2}r_ix^i$ and $q=\sum_{i\le d_1-d_2}q_ix^i$. Then, there is a circuit of size $O(d_1d_2)$, whose inputs are all $h_i,g_i$ and outputs are all $r_i,q_i$.
\paragraph{Complete Appendix E proof (Lemma 70).} We shall denote the desired circuit by $C_{d_1,d_2}$. So we want:
\begin{multline*}C_{d_1,d_2}(g_0,g_1,\ldots,g_{d_1},h_0,h_1,\ldots,h_{d_2})\\=(r_1,r_2,\ldots,r_{d_2-1},q_0,q_1,\ldots,q_{d_1-d_2}).\end{multline*}
If $d_1<d_2$, we know that $q=0$. Hence:
\[C_{d_2-1,d_2}(g_0,g_1,\ldots,g_{d_1-1},h_0,h_1,\ldots,h_{d_1})=(g_1,g_2,\ldots,g_{d_1-1}).\]
If $d_1>d_2$, we perform a long division step:
\begin{align*}g&\leftarrow g-hx^{d_1-d_2}g_{d_1}\\&=\sum_{i\le d_1-d_2-1}g_ix^i+\sum_{i\ge d_1-d_2}^{d_1-1}(g_i-h_{i-(d_1-d_2)}g_{d_1})x^i.\end{align*}
Note that, we can set $q_{d_1-d_2}=g_{d_1}$. Define:
\[\mathbf g\mathrel{\overset{\mathrm{def}}{=}}(g_0,g_1,\ldots,g_{d_1-d_2-1},g_{d_1-d_2}-h_0g_{d_1},\ldots,g_{d_1-1}-h_{d_2-1}g_{d_1}).\]
Then we have:
\begin{multline*}C_{d_1,d_2}(g_0,g_1,\ldots,g_{d_1},h_0,h_1,\ldots,h_{d_2})\\=(C_{d_1-1,d_2}(\mathbf g,h_0,h_1,\ldots,h_{d_2}),g_{d_1}).\tag{8}\end{multline*}
Hence if $S(d_1,d_2)$ is the size of $C_{d_1,d_2}$ then Equation (8) implies that $S(d_1,d_2)=S(d_1-1,d_2)+2d_2$ and $S(d_2-1,d_2)=2d_2-1$. Therefore $S(d_1,d_2)\le2d_1d_2$.\qed
\section*{Complete author monic transformation (Appendix C)}
Given any polynomial $p(\mathbf x)$ in variables $\mathbf x=(x_1,\ldots,x_n)$, there is a standard trick to make it monic in $x_n$ by applying a linear transformation on the variables: for $\alpha=(\alpha_1,\ldots,\alpha_{n-1})\in\mathbb F^{n-1}$, let $\tau_\alpha:x_i\mapsto\alpha_ix_n+x_i$, for $i\in[n-1]$, and $x_n\mapsto x_n$. Note that $\deg(\tau_\alpha(p))\le\deg(p)$ [it may decrease because of non-trivial cancellations]. It is easy to see that $\tau_\alpha$ is an invertible map. We show that $\tau_\alpha(p)$ is monic in $x_n$, for a random transformation $\tau_\alpha$ i.e. when $\alpha\in\mathbb F^{n-1}$ chosen randomly. In fact, we show that this map can simultaneously make polynomials monic given that the field $\mathbb F$ is sufficiently large.
\paragraph{Lemma 68 (Monic Transformation).} Let $p_1(\mathbf x),\ldots,p_m(\mathbf x)$ be $m$-many polynomial of degree $d$. Let $S\subseteq\mathbb F$ be a finite set. For $\alpha\in S^{n-1}$, picked independently and uniformly at random,
\[\Pr\left[\bigwedge_{i=1}^m\tau_\alpha(p_i(\mathbf x))\text{ is monic in }x_n\right]\ge1-\frac{dm}{|S|}.\]
\begin{proof}Consider the terms of degree $d$ of a non-zero polynomial $p\in\mathbb F[\mathbf x]$. Define the set
\[T:=\{\beta=(\beta_1,\ldots,\beta_n):|\beta|_0=\sum_i\beta_i=d,\text{ and }\operatorname{coef}_{\mathbf x^\beta}(p)\ne0\}.\]
We also denote $\beta'=(\beta_1,\ldots,\beta_{n-1})$, the first $n-1$-coordinates of $\beta$, and similarly $\mathbf x'=(x_1,\ldots,x_{n-1})$. Note that $\tau_\alpha(\mathbf x^\beta)=\alpha^{\beta'}x_n^d+(\text{lower terms in }x_n)$. Observe that the homogeneous component of degree $d$ in $\tau_\alpha(p)$ can be written as $a_{d,p}(\mathbf x)=\sum_{\beta\in T}c_\beta\tau_\alpha(\mathbf x^\beta)$, for some constants $c_\beta$. Trivially, $a_{d,p}$ is a nonzero polynomial, and moreover,
\[a_{d,p}(\alpha)=\left(\sum_{\beta\in T}c_\beta\alpha^{\beta'}\right)x_n^d+(\text{lower terms in }x_n).\]
In order to make $\tau_\alpha(p)$ monic in $x_n$, we want $(\sum_{\beta\in T}c_\beta\alpha^{\beta'})\ne0$. So, define, another polynomial $b_{d,p}(\mathbf x')=(\sum_{\beta\in T}c_\beta\mathbf x'^{\beta'})$. It can have degree atmost $d$. As we want each $\tau_\alpha(p)$ monic where $p=p_m(\mathbf x)$, it suffices to find $\alpha$ such that $\prod_{i\in[m]}b_{d,p_i}(\alpha)\ne0$. Note that $\deg(\prod_{i\in[m]}b_{d,p_i}(\mathbf x))\le d\cdot m$. Thus, when we pick $\alpha$ at random, the probability that $\prod_{i\in[m]}b_{d,p_i}(\alpha)=0$, is at most $\le dm/|S|$, from Lemma 64. Hence, the conclusion follows.\end{proof}

\section*{Complete author Section 3.2 main construction}
Now we prove the following Lemma 17 which shows that both $\mathrm{div}$ and $\mathrm{mod}$ have low complexity when the divisor has low-degree and monic (in fact, constant leading-coefficient suffices).
\paragraph{Lemma 17 (Main Lemma).} Let the polynomials $g,h\in\mathbb F[\mathbf x]$ such that $h$ is monic in $x_n$, $L(g)=s_1$, $L(h)=s_2$, and $\deg_{x_n}(h)=d$. Then, both $L(g\mathbin{\mathrm{div}}h),L(g\mathbin{\mathrm{mod}}h)\le O((s_1+s_2)d^2)$.
\begin{proof}Suppose $C$ is a circuit of size $s_2$ which computes $h$ and $C_g$ is the circuit of size $s_1$ which computes $g$. By using Lemma 61, there is a circuit of size $O(s_2d^2)$, which computes $h_0,\ldots,h_{d-1}$. Now, we split every gate $F$ in $C$ into $d+1$ gates $F_0,F_1,\ldots,F_d$. Suppose, the gate $F$ is computing a polynomial $P_F$. Let $P_F\mathbin{\mathrm{mod}}h=\sum_{i<d}p_ix_n^i$. Then we want the property that $F_i$ computes $p_i$ for $i<d$. And if $i=d$ then $F_i$ computes $P_F\mathbin{\mathrm{div}}h$.
Suppose $F$ is a $+$ gate in $C$ with children gates computing the polynomials $a$ and $b$. Again express $a\mathbin{\mathrm{mod}}h=\sum_{i<d}a_ix_n^i$ and $b\mathbin{\mathrm{mod}}h=\sum_{i<d}b_ix_n^i$. It is clear that
\[(a+b)\mathbin{\mathrm{mod}}h=a\mathbin{\mathrm{mod}}h+b\mathbin{\mathrm{mod}}h.\]
Therefore $p_i=a_i+b_i$. It is also clear that $P_F\mathbin{\mathrm{div}}h=a\mathbin{\mathrm{div}}h+b\mathbin{\mathrm{div}}h$.
Suppose $F$ is a $\times$ gate in $C$ with children gates computing the polynomials $a$ and $b$. Again express $a\mathbin{\mathrm{mod}}h=\sum_{i<d}a_ix_n^i$ and $b\mathbin{\mathrm{mod}}h=\sum_{i<d}b_ix_n^i$. It is clear that:
\[(a\cdot b)\mathbin{\mathrm{mod}}h=(a\mathbin{\mathrm{mod}}h\cdot b\mathbin{\mathrm{mod}}h)\mathbin{\mathrm{mod}}h.\]
For $\mathrm{div}$, we have that:
\begin{align*}P_F\mathbin{\mathrm{div}}h&=a\mathbin{\mathrm{div}}h\cdot b\mathbin{\mathrm{div}}h\cdot h+b\mathbin{\mathrm{div}}h\cdot a\mathbin{\mathrm{mod}}h\\&\quad+a\mathbin{\mathrm{div}}h\cdot b\mathbin{\mathrm{mod}}h+(a\mathbin{\mathrm{mod}}h\cdot b\mathbin{\mathrm{mod}}h)\mathbin{\mathrm{div}}h.\end{align*}
We have already computed $a\mathbin{\mathrm{mod}}h,b\mathbin{\mathrm{mod}}h,a\mathbin{\mathrm{div}}h,b\mathbin{\mathrm{div}}h$. So, we only need to compute $(a\mathbin{\mathrm{mod}}h\cdot b\mathbin{\mathrm{mod}}h)\mathbin{\mathrm{mod}}h$ and $(a\mathbin{\mathrm{mod}}h\cdot b\mathbin{\mathrm{mod}}h)\mathbin{\mathrm{div}}h$. Since we have already computed $a_i,b_i$ for all $i<d$, by using Lemma 16, we can compute all the $p_i$ and $(a\mathbin{\mathrm{mod}}h\cdot b\mathbin{\mathrm{mod}}h)\mathbin{\mathrm{div}}h$ in $O((2d-2)d)=O(d^2)$ many gates. Therefore the new circuit has $O(s_1d^2)$ has many gates. Also we used $O(s_2d^2)$ gates to computes $h_0,\ldots,h_{d-1}$. Hence,
\[L(g\mathbin{\mathrm{div}}h)=O((s_1+s_2)d^2),\qquad L(g\mathbin{\mathrm{mod}}h)=O((s_1+s_2)d^2).\]\end{proof}
The following theorem settles Problem 1, when the divisor has small degree (proving Theorem 4).
\paragraph{Proof of Theorem 18.} The above Lemma 17 shows that when $h$ is monic in $x_n$, the upper bound holds. Let $\tau:\mathbb F[\mathbf x]\longrightarrow\mathbb F[\mathbf x]$, be an invertible monic transformation (sends $x_i\mapsto\alpha_i\cdot x_n+x_i$, where $\alpha_i\in\mathbb F$) s.t. $\tau(h)$ is monic wrt $x_n$, such transformation exists (Lemma 68). Note that, $L(\tau(g))\le s+n=O(s_1)$, and $L(\tau(h))\le s_2+n=O(s_2)$. Moreover, as $\tau$ is degree-preserving, $\deg_{x_n}(\tau(h))=d$.
So, apply Lemma 17 to conclude that $\tau(f)=\tau(g)\mathbin{\mathrm{div}}\tau(h)$, has a circuit of size $O((s_1+s_2)d^2)$. We apply $\tau^{-1}$ again (which is just a additive $n$-blowup) to finally deduce that
\[L(f)\le O((s_1+s_2)d^2).\]\qed
\paragraph{Remark 19.} This proof holds when one replaces $L$ by $\tau$, i.e. the constant-free circuit complexity (for definition, see Section 5.2). Note that, neither $\mathrm{div}$ nor $\mathrm{mod}$ introduce any new constant in the process. Moreover, one can choose the $\alpha_i$ to be explicit and $\mathrm{poly}(\log d)$-computable so that $\tau$ is a monic invertible map. This establishes the claim.
\section*{Complete author Section 3.3 border proof}
\paragraph{Proof of Theorem 22.} By definition, there exists $G,H\in\mathbb F(\epsilon)[\mathbf x]$, of size at most $s_1$ and $s_2$, respectively, such that $G:=g+\epsilon\cdot\widetilde g(\mathbf x,\epsilon)$, and $H:=h+\epsilon\cdot\overline h(\mathbf x,\epsilon)$, where $\widetilde g,\overline h\in\mathbb F[\epsilon,\mathbf x]$. We note that, $\deg_{\mathbf x}(H)$ can be larger than $d$. However, using Lemma 62, we know that $L_{\mathbb F(\epsilon)}(\mathrm{Hom}_{\le d}H)\le O(s_2d^2)=:s'_2$.
We denote $\widetilde H:=\mathrm{Hom}_{\le d}H$. It is important to observe that $\widetilde H|_{\epsilon=0}=h$. By definition, there exists $m$ (could be $\exp(s'_2)$) such that
\[\widetilde H:=h+\epsilon\cdot\widetilde h(\mathbf x,\epsilon)=h+\sum_{j\in[m]}\epsilon^jh_j(\mathbf x),\qquad h_j\in\mathbb F[\mathbf x].\]
Let $\tau:\mathbb F[\mathbf x]\longrightarrow\mathbb F[\mathbf x]$, be an invertible monic transformation (sends $x_i\mapsto\alpha_i\cdot x_n+x_i$, where $\alpha_i\in\mathbb F$) s.t. $\tau(h)$ and each $\tau(h_j)$, for $j\in[m]$ is monic wrt $x_n$; such transformation exists (Lemma 68). Note that, $L_{\mathbb F(\epsilon)}(\tau(G))\le O(s_1)$ and $L_{\mathbb F(\epsilon)}(\tau(\widetilde H))\le O(s'_2)$. Further, $\deg_{\mathbf x}(\tau(\widetilde H))=d$, as $\tau$ is a degree-preserving map. We also have the following identities:
\[\tau(\widetilde H)=\tau(h)+\epsilon\cdot\tau(\widetilde h),\qquad\tau(G)=\tau(f)\cdot\tau(h)+\epsilon\cdot\tau(\widetilde g).\]
By assumption, the leading coefficient of $x_n$ in $\tau(\widetilde H)$ (call it $\alpha$) is in $\mathbb F[\epsilon]$ (in fact, $\alpha\not\equiv0\bmod\epsilon$). This basically makes $\tau(\widetilde H)$ a monic polynomial over $\mathbb F(\epsilon)[\mathbf x]$. Therefore, $\mathrm{div}\,\tau(\widetilde H)$ and $\mathrm{mod}\,\tau(\widetilde H)$ now make sense over $\mathbb F(\epsilon)[\mathbf x]$. By simple division, we have
\begin{multline*}\tau(G)\mathbin{\mathrm{div}}\tau(\widetilde H)=\tau(f)\\+\epsilon\cdot(\tau(\widetilde g)-\tau(f)\cdot\tau(\widetilde h))\mathbin{\mathrm{div}}\tau(\widetilde H).\tag{1}\end{multline*}
Note that, Lemma 17 implies $L_{\mathbb F(\epsilon)}(\tau(G)\mathbin{\mathrm{div}}\tau(\widetilde H))=O((s_1+s'_2)d^2)$. By definition of $\underline L$ and Equation (1), it is trivial to conclude that $\underline L(\tau(f))\le O((s_1+s'_2)d^2)=O(s_1d^2+s_2d^4)$. As $\tau$ is invertible, we can get back $f$ by applying $\tau^{-1}$ (incurring $n$-additive blowup). This finally shows
\[\underline L(f)\le O(s_1d^2+s_2d^4).\]\qed

\section*{Formal author rational-series truncation result}
\paragraph{Theorem 26 (Truncation of low-degree rational function).} Suppose, $g$ and $h$ are two univariate polynomials in $\mathbb F[x]$ such that $\deg(g)\le d$, $\deg(h)=d_h$, and $g$ can be computed a circuit of size $s$. Let, $g/h\in\mathbb F[[x]]$. Then, truncation of $g/h$ upto degree-$d$ can be computed by a circuit of size $\mathrm{poly}(s,d_h,\log d)$.

\section*{Complete author Section 5.1 proof and required mathematical tools}
A formal power series is a generalization of a polynomial, where the number of terms can be infinite. Formally, $A=\sum_{i\ge0}A_ix^i$ with $A_i\in\mathbb F$, is a power series in the power series ring $\mathbb F[[x]]$. We define the degree $d$ truncation $\mathrm{trunc}(A,d)$ of $A$ to be $\mathrm{trunc}(A,d):=\sum_{0\le i\le d}A_ix^i$. So, $\mathrm{trunc}(A,d)$ is always a polynomial of degree at most $d$.
We show that complexity of truncation of rational functions where the degrees are small, has low complexity. For simplicity, we work with $\mathbb F=\overline{\mathbb F}$, an algebraically closed field. We first recall the following folklore decomposition.
\paragraph{Lemma 24 (Partial fraction decomposition).} Let $g(x)/h(x)$ be a rational function with $\deg(g)<\deg(h)$. If $h(x)=\prod_{i\in[k]}(x-a_i)^{d_i}$ is the factorization of $h(x)$ over $\mathbb F[x]$, then, there exist $b_{ij}\in\mathbb F$ s.t.:
\[g(x)/h(x)=\sum_{i\in[k]}\sum_{j\in[d_i]}b_{ij}/(x-a_i)^j.\]
\paragraph{Lemma 25.} For any non zero $a\in\mathbb F$, we have $L(\mathrm{trunc}(1/(x-a),d))=O(\log d)$.
\begin{proof}This follows from the inverse identity: $1/(a-x)=1/a\sum_{i\ge0}(x/a)^i$ and the fact that $L(\sum_{0\le i\le d}(x/a)^i)=O(\log d)$ (By using Remark 15).\end{proof}
\paragraph{Lemma 65 (Power series valuation).} Let $g,h\in\mathbb F[x]$ such that $g/h\in\mathbb F[[x]]$. Let $m$ (respec. $n$) be the highest power dividing $g$ (respec. $h$) i.e. $x^m\mid g$ and $x^{m+1}\nmid g$ (respec. for $h$). Then, $m\ge n$.
\begin{proof}Suppose, $m<n$. Note that, there exists $0\ne\alpha\in\mathbb F$, such that $h=\alpha x^n(1+x\widetilde h)$, for some $\widetilde h\in\mathbb F[x]$. Similarly, let $g=\beta x^m(1+x\widetilde g)$, for some $\widetilde g\in\mathbb F[x]$ and $\beta\in\mathbb F$. Thus,
\begin{align*}\frac gh&=\frac\beta\alpha x^{m-n}\frac{1+x\widetilde g}{1+x\widetilde h}\\&=\frac\beta\alpha x^{m-n}(1+x\widetilde g)(1+x\widetilde h+(x\widetilde h)^2+\cdots)\\&\notin\mathbb F[[x]],\text{ a contradiction}.\end{align*}\end{proof}
\paragraph{Lemma 66 (General Leibniz rule).} If $f$ and $g$ are $k$-time differentiable functions, then $(fg)^{(k)}=\sum_{i=0}^k\binom ki f^{(k-i)}g^{(i)}$.
\paragraph{Lemma 63.} Let $f$ be a polynomial $\mathbb F[\mathbf x]$, computed by a size $s$ circuit $C$. Then, there exists a circuit $C'$ of size $O(sm^2)$ which computes $f,f^{(1)},f^{(2)},\ldots,f^{(m)}$.
\begin{proof}We split every $G$ gate in $C$ to $n+1$ gates $G_0,\ldots,G_m$ in $C'$. The property we want is that if the gate $G$ is computing the polynomial $g$ in $C$ then $G_k$ computes the polynomial $g^{(k)}$ in $C'$. Suppose $G$ is a $+$ gate in $C$ with children gates computing the polynomials $g_1$ and $g_2$. Now we know that $g^{(k)}=g_1^{(k)}+g_2^{(k)}$. Thus we can easily propagate the derivatives on addition/subtraction gates. If $G$ is a $\times$ gate then using Lemma 66, we know that:
\[(g_1g_2)^{(k)}=\sum_{i=0}^k\binom ki g_1^{(k-i)}g_2^{(i)}.\]
Thus we can computes $g,g^{(1)},g^{(2)},\ldots,g^{(m)}$ using additional $O(m^2)$ gates. Therefore $C'$ has $O(sm^2)$ gates.\end{proof}
\paragraph{Proof of Theorem 26.} The main idea is to use Lemma 24 and the low complexity of the truncation of inverse identity (Lemma 25). However, the given polynomial $h$ may be divisible by $x$ (i.e. $h(0)=0$). In that case, let $m$ be the highest power of $x$ such that $x^m\mid h$ (i.e. $x^{m+1}\nmid h$). Note that, as $g/h\in\mathbb F[[x]]$, $x^m\mid g$ as well (Lemma 65). As $\deg(h)\le d_h$, thus $m\le d_h$.
By using Theorem 18, we know that $g_1:=g/x^m$ has a cicuit of size $O((s+\log d_h)d_h^2)=:s_1$. Trivially, $h_1:=h/x^m$ has degree $\le d_h$, and $g/h=g_1/h_1$. Denote, $g_2:=g_1\mathbin{\mathrm{mod}}h_1$. Obviously, $\deg(g_2)<\deg(h_1)$ and $g_1/h_1=g_1\mathbin{\mathrm{div}}h_1+g_2/h_1$. Invoking Theorem 13, one concludes that $L(g_1\mathbin{\mathrm{div}}h_1)=O(s_1d_h)$. Therefore, $L(\mathrm{trunc}(g_1\mathbin{\mathrm{div}}h_1,d))=O(s_1d_h)$, as $\deg(g_1)<d$.
Let $h_1$ factors over $\mathbb F[x]$ as $h_1:=\prod_{i\in[k]}(x-a_i)^{d_i}$. Trivially, $\sum d_i\le d_h$. By using Lemma 24 on $g_2/h_1$, we know that there are constants $a_i,b_{ij}\in\mathbb F$ such that:
\[g_2(x)/h_1(x)=\sum_{i\in[k]}\sum_{j\in[d_i]}b_{ij}/(x-a_i)^j.\]
Note that, for any $a\in\mathbb F$ and $t\in\mathbb N$, $d^t/dx^t(1/(x-a))=(-1)^tt!\cdot1/(x-a)^{t+1}$, and thus,
\begin{multline*}\mathrm{trunc}(1/(x-a)^{t+1},d)\\=\frac{(-1)^t}{t!}\frac{d^t}{dx^t}(\mathrm{trunc}(1/(x-a),d))+\sum_{i=d-t+1}^d\gamma_ix^i,\qquad\gamma_i\in\mathbb F.\end{multline*}
Using the above identity and Lemma 63, we can show that
\[L\left(\mathrm{trunc}\left(\sum_{j\in[d_i]}b_{ij}/(x-a_i)^j,d\right)\right)=O(\log d\cdot d_i^2).\]
To show this, note that $L(\mathrm{trunc}(1/(x-a_i),d))=O(\log d)$, and using Lemma 63, we compute all its derivative till the $d_i$-th one which has a circuit of size $O(\log d\cdot d_i^2)$. Using the above identity, we can add $d_i-1$ many monomials of the form $cx^\ell$ with $d-d_i+2\le\ell\le d$ (each monomial has trivial size of $O(\log d)$) to the circuit to obtain a circuit for $\mathrm{trunc}(\sum_{j\in[d_i]}b_{ij}/(x-a_i)^j,d)$, which still has size $O(\log d\cdot d_i^2)$. Thus, doing it for each $a_i$ for $i\in[k]$, one obtains that
\begin{align*}L(\mathrm{trunc}(g(x)/h(x),d))&=L(\mathrm{trunc}(g_1(x)/h_1(x),d))\\&=L(g_1\mathbin{\mathrm{div}}h_1,d)+L(g_2/h_1,d)\\&=O(s_1d_h)+L\left(\mathrm{trunc}\left(\sum_{i\in[k]}\sum_{j\in[d_i]}b_{ij}/(x-a_i)^j,d\right)\right)\\&=O((s+\log d_h)d_h^3)+O\left(\log d\cdot\sum_{i\in[k]}d_i^2\right)\\&=O(sd_h^3\log d).\end{align*}\qed
\paragraph{Remark 27.} Eventually, we can replace $g\in\mathbb F[[x]]$ with the given complexity $\mathrm{trunc}(g,d)=s$ and show that the exact same proof as above, works.

\begin{thebibliography}{99}
\bibitem{ref1} Alexander Alder. Grenzrang und Grenzkomplexität aus algebraischer und topologischer Sicht. PhD thesis, Zentralstelle der Studentenschaft, 1984.
\bibitem{ref2} Eric Allender, Peter Bürgisser, Johan Kjeldgaard-Pedersen, and Peter Miltersen. On the Complexity of Numerical Analysis. SIAM Journal on Computing, 38, January 2006. Preliminary version in the 21st Annual IEEE Conference on Computational Complexity (CCC’06). doi: 10.1109/CCC.2006.30.
\bibitem{ref3} Robert Andrews. Algebraic Hardness Versus Randomness in Low Characteristic. In 35th Computational Complexity Conference (CCC 2020), volume 169 of Leibniz International Proceedings in Informatics (LIPIcs), pages 37:1–37:32. Schloss Dagstuhl–Leibniz-Zentrum für Informatik, 2020. doi:10.4230/LIPIcs.CCC.2020.37.
\bibitem{ref4} Jean Berstel and Christophe Reutenauer. Rational Series and Their Languages. Springer- Verlag, Berlin, Heidelberg, 1988. URL: https://dl.acm.org/doi/book/10.5555/52107.
\bibitem{ref5} Markus Bläser and Gorav Jindal. On the Complexity of Symmetric Polynomials. In 10th Innovations in Theoretical Computer Science Conference (ITCS’19), volume 124 of LIPIcs, pages 47:1–47:14. Schloss Dagstuhl - Leibniz-Zentrum für Informatik, 2019. doi:10.4230/ LIPIcs.ITCS.2019.47.
\bibitem{ref6} Alin Bostan, Gilles Christol, and Philippe Dumas. Fast computation of the Nth term of an algebraic series over a finite prime field. In Proceedings of the ACM on International Symposium on Symbolic and Algebraic Computation (ISSAC’16), pages 119–126, 2016. doi: 10.1145/2930889.2930904.
\bibitem{ref7} Karl Bringmann, Christian Ikenmeyer, and Jeroen Zuiddam. On Algebraic Branching Programs of Small Width. J. ACM, 65(5):1–29, 2018. (Preliminary version in the 32nd Computational Complexity Conference (CCC’17). doi:10.1145/3209663.
\bibitem{ref8} Peter Bürgisser. The complexity of factors of multivariate polynomials. Foundations of Computational Mathematics, 4(4):369–396, 2004. arXiv:1812.06828.
\bibitem{ref9} Peter Bürgisser. On defining integers and proving arithmetic circuit lower bounds. Computa- tional Complexity, 18(1):81–103, 2009. doi:10.1007/s00037-009-0260-x.
\bibitem{ref10} Peter Bürgisser, Michael Clausen, and Amin Shokrollahi. Algebraic complexity theory, volume 315. Springer Science \& Business Media, 2013.
\bibitem{ref11} David V Chudnovsky and Gregory V Chudnovsky. On expansion of algebraic functions in power and Puiseux series, I. Journal of Complexity, 2(4):271–294, 1986. URL: https: //www.sciencedirect.com/science/article/pii/0885064X86900063.
\bibitem{ref12} Michael Coons. The Transcendence of Series Related to Stern’s Diatomic Sequence. Interna- tional Journal of Number Theory, 06, November 2011. doi:10.1142/S1793042110002958.
\bibitem{ref13} Michael Coons and Peter Borwein. Transcendence of power series for some number theoretic functions. Proceedings of the American Mathematical Society, 137, July 2008. doi:10.1090/ S0002-9939-08-09737-2.
\bibitem{ref14} Wellington De Melo and Benar Fux Svaiter. The cost of computing integers. Proceedings- American Mathematical Society, 124:1377–1378, 1996. URL: https://www.ams.org/journals/ proc/1996-124-05/S0002-9939-96-03173-5/S0002-9939-96-03173-5.pdf.
\bibitem{ref15} Richard A. Demillo and Richard J. Lipton. A probabilistic remark on algebraic program testing. Information Processing Letters, 7(4):193–195, 1978. URL: https://www.sciencedirect.com/ science/article/abs/pii/0020019078900674.
\bibitem{ref16} Pranjal Dutta. Real tau-Conjecture for sum-of-squares: A unified approach to lower bound and derandomization. In 16th International Computer Science Symposium in Russia (CSR 2021), 2021. URL: https://drive.google.com/file/d/1X8eo9GM4SCNsC2vWjPbUwMX0vff5i2k3/ view.
\bibitem{ref17} Pranjal Dutta, Nitin Saxena, and Amit Sinhababu. Discovering the roots: Uniform closure results for algebraic classes under factoring. In Proceedings of the 50th Annual ACM SIGACT Symposium on Theory of Computing, pages 1152–1165, 2018. URL: https://www.cse.iitk. ac.in/users/nitin/papers/factor-closure.pdf.
\bibitem{ref18} Pranjal Dutta, Nitin Saxena, and Thomas Thierauf. A Largish Sum-Of-Squares Implies Circuit Hardness and Derandomization. In 12th Innovations in Theoretical Computer Science Conference (ITCS 2021), volume 185 of Leibniz International Proceedings in Informatics (LIPIcs), pages 23:1–23:21. Schloss Dagstuhl–Leibniz-Zentrum für Informatik, 2021. URL: 10.4230/LIPIcs.ITCS.2021.23.
\bibitem{ref19} Joshua A Grochow et al. Complexity in ideals of polynomials: questions on algebraic complexity of circuits and proofs. Bulletin of EATCS, 2(130), 2020. URL: http://bulletin.eatcs.org/ index.php/beatcs/article/view/607.
\bibitem{ref20} Joshua A. Grochow, Ketan D. Mulmuley, and Youming Qiao. Boundaries of VP and VNP. In 43rd International Colloquium on Automata, Languages, and Programming (ICALP 2016), volume 55, pages 34:1–34:14, 2016. URL: https://core.ac.uk/download/pdf/62922137.pdf.
\bibitem{ref21} Gorav Jindal. On approximate polynomial identity testing and real root finding. PhD thesis, Saarland University, 2019. doi:10.22028/D291-29880.
\bibitem{ref22} Erich Kaltofen. Uniform closure properties of p-computable functions. In Proceedings of the eighteenth annual ACM symposium on Theory of computing, pages 330–337, 1986. doi: 10.1145/12130.12163.
\bibitem{ref23} Erich Kaltofen. Single-factor Hensel lifting and its application to the straight-line complexity of certain polynomials. In Proceedings of the 19th annual ACM symposium on Theory of computing (STOC’87), pages 443–452, 1987. doi:10.1145/28395.28443.
\bibitem{ref24} Pascal Koiran. Valiant’s model and the cost of computing integers. computational complexity, 13(3):131–146, 2005.
\bibitem{ref25} Pascal Koiran. Shallow circuits with high-powered inputs. Innovations in Computer Sci- ence (ICS), 2011. URL: https://hal-ens-lyon.archives-ouvertes.fr/ensl-00477023v4/ document.
\bibitem{ref26} Pascal Koiran and Sylvain Perifel. Interpolation in Valiant’s theory. Computational Complexity, 20(1):1–20, 2011. doi:10.1007/s00037-011-0002-8.
\bibitem{ref27} Pascal Koiran, Natacha Portier, Sébastien Tavenas, and Stéphan Thomassé. A $\tau$ -Conjecture for Newton Polygons. Foundations of computational mathematics, 15(1):185–197, 2015. doi: 10.1007/s10208-014-9216-x.
\bibitem{ref28} FV Kuhlmann. On convergent power series, 1996. URL: https://www.mathi.uni-heidelberg. de/~roquette/KONVPOTREIHEN.pdf.
\bibitem{ref29} Hsiang Kung and Joseph Traub. All Algebraic Functions Can Be Computed Fast. J. ACM,
\bibitem{ref30} Dick Lipton and Ken Regan. Factoring and factorials, February 2009. URL: https://rjlipton. wordpress.com/2009/02/23/factoring-and-factorials/.
\bibitem{ref31} Richard J Lipton. Polynomials with 0-1 coefficients that are hard to evaluate. SIAM Journal on Computing, 7(1):61–69, 1978. Preliminary version in the 16th Annual Symposium on Foundations of Computer Science (FOCS 1975). URL: https://epubs.siam.org/doi/abs/ 10.1137/0207004?journalCode=smjcat.
\bibitem{ref32} Richard J Lipton. Straight-line complexity and integer factorization. In International Algorithmic Number Theory Symposium (ANTS 94), pages 71–79. Springer, 1994. doi: 10.1007/3-540-58691-1\_45.
\bibitem{ref33} Richard J Lipton and Larry J Stockmeyer. Evaluation of polynomials with super- preconditioning. Journal of Computer and System Sciences, 16(2):124–139, 1978. URL: https://www.sciencedirect.com/science/article/pii/0022000078900417.
\bibitem{ref34} Meena Mahajan. Algebraic Complexity Classes. In Perspectives in Computational Complexity, pages 51–75. Springer, 2014. doi:10.1007/978-3-319-05446-9\_4.
\bibitem{ref35} Carlos Moreira. On asymptotic estimates for arithmetic cost functions. Proceedings of the American Mathematical Society, 125(2):347–353, 1997. URL: https://www.jstor.org/ stable/2161660.
\bibitem{ref36} Ketan D Mulmuley. The GCT program toward the P vs. NP problem. Communications of the ACM, 55(6):98–107, 2012. doi:10.1145/2184319.2184341.
\bibitem{ref37} Danny Nguyen and Igor Pak. Complexity of short generating functions. In Forum of Mathematics, Sigma, volume 6. Cambridge University Press, 2018. arXiv:1702.08660.
\bibitem{ref38} Øystein Ore. Über höhere kongruenzen. Norsk Mat. Forenings Skrifter, 1(7):15, 1922.
\bibitem{ref39} Igor Pak. Complexity problems in enumerative combinatorics. In Proceedings of the Inter- national Congress of Mathematicians – Rio de Janeiro 2018. Vol. IV. Invited lectures, pages 3153–3180. World Sci. Publ., Hackensack, NJ, 2018. doi:10.1142/9789813272880\_0176.
\bibitem{ref40} David A Plaisted. New NP-hard and NP-complete polynomial and integer divisibility problems. Theoretical Computer Science, 31(1-2):125–138, 1984. Preliminary in the 17th Annual Symposium on Foundations of Computer Science (FOCS 1976). URL: https: //www.sciencedirect.com/science/article/pii/0304397584901300.
\bibitem{ref41} Jacob T Schwartz. Fast probabilistic algorithms for verification of polynomial identities. Journal of the ACM (JACM), 27(4):701–717, 1980. doi:10.1145/322217.322225.
\bibitem{ref42} Adi Shamir. Factoring numbers in O (logn) arithmetic steps. Information Processing Let- ters, 8(1):28–31, 1979. URL: https://www.sciencedirect.com/science/article/abs/pii/ 0020019079900875.
\bibitem{ref43} Amir Shpilka and Amir Yehudayoff. Arithmetic Circuits: A survey of recent results and open questions. Foundations and Trends® in Theoretical Computer Science, 5(3–4):207–388, 2010. doi:10.1561/0400000039.
\bibitem{ref44} Michael Shub and Steve Smale. On the intractability of Hilbert’s Nullstellensatz and an algebraic version of “NP $\neq$ P?”. Duke Math. J., 81(1):47–54 (1996), 1995. A celebration of John F. Nash, Jr. doi:10.1215/S0012-7094-95-08105-8.
\bibitem{ref45} Volker Strassen. Vermeidung von divisionen. Journal für die reine und angewandte Mathematik, 264:184–202, 1973.
\bibitem{ref46} L Valiant. Reducibility by algebraic projections in: Logic and algorithmic. In Symposium in honour of Ernst Specker, pages 365–380, 1982.
\bibitem{ref47} Leslie G Valiant. Completeness classes in algebra. In Proceedings of the 11th Annual ACM symposium on Theory of computing, pages 249–261. ACM, 1979. doi:10.1145/800135.804419.
\bibitem{ref48} Joachim Von Zur Gathen and Jürgen Gerhard. Modern computer algebra. Cambridge university press, 2013.
\bibitem{ref49} Wact. Some Accessible Open Problems. Workshop on Algebraic Complexity Theory (WACT 2016).   URL: https://www.cs.tau.ac.il/~shpilka/wact2016/concreteOpenProblems/ openprobs.pdf.
\bibitem{ref50} Klaus W Wagner. The complexity of combinatorial problems with succinct input representation. Acta informatica, 23(3):325–356, 1986. doi:10.1007/BF00289117.
\bibitem{ref51} Richard Zippel. Probabilistic Algorithms for Sparse Polynomials. In Proceedings of the International Symposium on Symbolic and Algebraic Computation, EUROSAM ’79, pages 216–226, 1979. doi:10.1007/3-540-09519-5\_73.
\bibitem{ref52} J.von zur Gathen and V. Strassen. Some polynomials that are hard to compute. Theoretical Computer Science, 11(3):331–335, 1980. URL: http://www.sciencedirect.com/science/ article/pii/0304397580900201.
\end{thebibliography}
\end{document}
